% TOP_PROBLEM: {"problemId": "problem.the-nc-versus-p-problem-for-efficient-parallel-computation", "canonicalTitle": "The NC versus P Problem for Efficient Parallel Computation", "releaseRank": 23, "primaryDomain": "theoretical_computer_science", "primaryDomainLabel": "Theoretical computer science", "displayStatus": "open", "releaseStatus": "open"}
% !TeX program = lualatex
% ATTEMPT_STATUS: unresolved
\documentclass[11pt]{article}
\usepackage[margin=25mm]{geometry}
\usepackage{fontspec}
\setmainfont{Latin Modern Roman}
\usepackage{amsmath,amssymb,mathtools}
\usepackage{unicode-math}
\setmathfont{Latin Modern Math}
\usepackage{xurl,hyperref}
\hypersetup{colorlinks=true,urlcolor=blue}
\setcounter{secnumdepth}{0}
\setlength{\parindent}{0pt}
\setlength{\parskip}{5pt}
\setlength{\emergencystretch}{3em}
\begin{document}
\section*{\texorpdfstring{The NC versus P Problem for Efficient Parallel Computation}{The NC versus P Problem for Efficient Parallel Computation}}\label{23}
\subsection{Definitions and mathematical statement}
For bounded-fan-in Boolean circuits, size is the number of gates and depth is the longest input-to-output path. In the standard uniform convention, $\mathsf{NC}=\bigcup_{k\ge1}\mathsf{NC}^k$, where $\mathsf{NC}^k$ has polynomial size and depth $O((\log n)^k)$. Uniformity requires an effective resource-bounded procedure generating the circuit family, rather than an arbitrary circuit for each length. The question is
\[
 \mathsf{NC}\stackrel{?}{=}\mathsf P.
\]
For precision, the conventional logspace-uniform interpretation is used; the catalogue does not choose among finer uniformity conventions.
\par\smallskip\textit{Formulation and concept references:} \href{https://arxiv.org/abs/2508.20735}{[S1]}.

\subsection{Short English statement}
Can every efficient sequential decision algorithm be replaced by a highly parallel one using only polylogarithmically many stages?
\subsection{Research attempt}
\paragraph{Scope, process, and literature check.}
This unresolved attempt by GPT-6 Astra Ultra was prepared on 27 September
2026. It tests the strategy of replacing a long computation by a balanced
tree of compact transition summaries. Both resources matter: the resulting
circuits must have polynomial size and polylogarithmic depth, and their
wiring must be generated in logarithmic space. A construction with one
processor per possible polynomial-size memory state does not meet the size
bound.

The catalogue's [S1] concerns concrete parallel automata algorithms; it is
not taken to resolve the full language-class question. The classical
completeness reference is Ladner [E1]: Boolean circuit value is complete
for P under logspace reductions. This gives a concrete universal target.
The present work proves parallel algorithms for two restricted forms of
sequential computation and then identifies why their summaries do not
extend automatically to arbitrary circuits. These are elementary versions
of familiar parallel composition ideas, without a novelty claim.

The easy inclusion is $\mathsf{NC}\subseteq\mathsf P$. Generate the
polynomial-size circuit using its uniformity machine and evaluate it gate
by gate in a topological order. Its size and encoding length are polynomial,
so this takes polynomial time. A standard halting logspace generator with
polynomial output can itself be simulated in polynomial time. Thus equality
would follow from a uniform polynomial-size, polylogarithmic-depth evaluator
for arbitrary polynomial-size Boolean circuits. Conversely, if NC equals P,
such an evaluator exists because circuit value lies in P.

\paragraph{Attack plan: summarize before balancing.}
Associativity permits balancing a chain of compositions without changing
its value. That alone does not bound the size of the objects being composed,
or the parallel cost of composition. We first use explicit transition tables,
where every bound can be counted. We then replace tables by affine maps,
which can describe transitions on exponentially many states with polynomially
many bits. The central question is whether arbitrary circuit computations
admit an equally effective summary language.

\paragraph{Lemma 23.1: explicit finite-state transitions.}
Suppose $m$ deterministic transitions $f_1,\ldots,f_m$ act on a state set
$[w]$. Given their tables and an initial state $a$, the final state
$f_m\circ\cdots\circ f_1(a)$ is computable by bounded-fan-in circuits with
\[
 \text{size }O(mw^3+mw^2\log(w+2)),\qquad
 \text{depth }O(\log(m+2)\log(w+2)).
\]
For standard padded table encodings, the circuit wiring is logspace-uniform
as a function of the dimensions. In particular the result is in uniform
NC when $m,w$ are polynomially bounded in the actual input length.

\emph{Proof.} Represent $f_i$ by a Boolean matrix $M_i$ with
$(M_i)_{u,v}=1$ exactly when $f_i(u)=v$. Form matrices for adjacent blocks
by Boolean matrix multiplication:
\[
                (BC)_{u,v}=\bigvee_{z=1}^w(B_{u,z}\wedge C_{z,v}).
\]
If $B$ and $C$ represent successive deterministic transitions, their product
represents their composition in that order. The formula also works for
arbitrary binary relations, so no unproved algebraic identity is needed.
Each entry takes $w$ AND gates and a balanced binary OR tree. A whole
matrix product has $O(w^3)$ gates and depth $O(\log(w+2))$.

Balance the $m$ matrices in a binary tree, padding with identities to the
next power of two if necessary. There are fewer than $2m$ leaves and
$O(m)$ internal products. The tree has $O(\log(m+2))$ levels. Converting
binary table entries to the matrices uses parallel equality comparisons
with $O(\log(w+2))$ gates per comparison. Reading the final row indexed by
$a$ and encoding its unique nonzero column has smaller polynomial cost
and at most $O(\log(w+2))$ additional depth.

For uniformity, a gate is described by its tree level, block number,
row, column, summation index, and position in an OR tree. Each field uses
$O(\log(mw+2))$ bits. The source-gate labels are obtained by fixed index
arithmetic on these fields. A generator loops through their polynomially
many possible values using logarithmic workspace and prints the gates and
wires. The equality-comparison and input-selection subcircuits have the
same type of explicit indexing.\hfill$\square$

When transition tables are the input, the bound is polynomial in their
encoding length. When a succinct circuit describes a transition on $b$
state bits, however, the table has $w=2^b$ states. Substitution into the
lemma gives exponentially many gates in $b$, and its stated depth contains
a factor $\log w=b$. For $b=\operatorname{poly}(n)$ this proves neither
required NC bound. Succinctness of the input transition cannot be confused
with explicit availability of its whole table.

\paragraph{Lemma 23.2: affine summaries avoid state enumeration.}
Let $m$ updates on $b$ bits have the form
$f_i(v)=A_i v+c_i$ over $\mathbb F_2$, with their coefficients explicitly
given. Their composition and its action on a given initial vector have
logspace-uniform circuits of polynomial size and depth
\[
                     O(\log(m+2)\log(b+2)).
\]
This remains useful when $b$ itself is polynomial in the input length.

\emph{Proof.} Augment $v$ by a final coordinate $1$. Then $f_i$ is represented
by the $(b+1)\times(b+1)$ matrix
\[
                   \widetilde A_i=
                   \begin{pmatrix}A_i&c_i\\0&1\end{pmatrix}.
\]
Composition is ordinary matrix multiplication over $\mathbb F_2$. Each
matrix entry is the XOR of $b+1$ products of bits. A binary XOR is a
constant-size AND/OR/NOT circuit, and balancing the XOR gives depth
$O(\log(b+2))$ and size $O(b+1)$ per entry. A balanced product tree
therefore has size $O(m(b+1)^3)$ and the stated depth. Multiplication by
the initial augmented vector adds one more such layer. Gate labels and
wiring are enumerated by the same logarithmic-space index scheme as in
Lemma 23.1, replacing OR by XOR.\hfill$\square$

The distinction between these lemmas is meaningful. The affine map can act
on $2^b$ states, yet its summary uses only $O(b^2)$ bits and remains in the
same format after composition. The exact formula
\[
 (A_2,c_2)\circ(A_1,c_1)
       =(A_2A_1,A_2c_1+c_2)
\]
shows both closure and the cost of implementing it. An existential claim
that each block has a short description would be weaker: NC also needs
the descriptions and their combination to be computed uniformly in shallow
circuits.

\paragraph{A nonlinear example that still admits short summaries.}
Affine maps do not capture all useful parallelism. In binary addition,
let $a_i,b_i$ be the input bits, let $g_i=a_i\wedge b_i$ and
$p_i=a_i\mathbin{\oplus}b_i$, and write the incoming carry as $c_i$.
The next carry is
\[
                         c_{i+1}=g_i\vee(p_i\wedge c_i).
\]
Although this describes a sequential dependence through all positions,
each block acts on its incoming carry by a map
$T_{g,p}(c)=g\vee(p\wedge c)$. These maps satisfy
\begin{equation}\label{a23:carry}
 T_{g_2,p_2}\circ T_{g_1,p_1}
  =T_{g_2\vee(p_2\wedge g_1),\,p_2\wedge p_1}.
\end{equation}
Distribute the second expression over the first to verify the identity.
Associativity follows from function composition, even when distinct pairs
represent the same constant function. Thus a block has only two summary bits
and combines in constant depth.

For each prefix, independently balance its sequence of pairs. This elementary
construction uses at most $O(n^2)$ gates overall and depth $O(\log(n+2))$;
sharing prefix computations can improve size but is unnecessary for NC.
Each carry and then each sum bit is computed in parallel. The positions
and balanced-tree gate labels are logspace-enumerable. Hence a linear chain
of dependencies in a particular sequential algorithm is not a lower bound
on the depth of all circuits computing its result.

\paragraph{Two failed extensions to arbitrary circuits.}
First, summarize a general block by its input-output truth table and apply
Lemma 23.1. A block with $b$ live Boolean values can require a table with
$2^b$ input rows, even if its gate description is short. An identity block
already makes this tabular format exponentially larger than its simple
symbolic description. This rejects the chosen representation bound, not the
possibility of other compact summaries. The affine and carry examples show
that suitable algebra can avoid tables; a general replacement has not been
found here.

Second, unfold all shared subcomputations into a formula and try to balance
the resulting tree. Unfolding need not have polynomial size. Let
$v_1=x$, $v_2=y$, and $v_i=v_{i-1}\wedge v_{i-2}$ for $i\ge3$.
There are only linearly many gates. The numbers of leaves in the mechanically
unfolded trees obey $\ell_i=\ell_{i-1}+\ell_{i-2}$, with
$\ell_1=\ell_2=1$, and grow exponentially. Nevertheless $v_i=x\wedge y$
for every $i\ge3$, by immediate induction. Thus the unfolded expression is
large even though a constant-depth circuit computes the function.
The blow-up disproves the efficiency of that preprocessing step, not NC
membership of the problem. Recognizing simplifications in this easy example
is not a general balancing algorithm for arbitrary shared circuits.

\paragraph{Connection to the full target and uniformity audit.}
Circuit value is in P by topological evaluation. For completeness, a fixed
deterministic machine running for $n^k$ steps can be represented by a
polynomial-size Boolean circuit: make gates for each local tape-cell update
at each time, with the state and head information encoded in the cells.
Gate positions use counters of $O(\log n)$ bits because $k$ is fixed.
The reduction prints the local wiring and initial input values using these
counters, yielding the logspace reduction underlying [E1].

Logspace reductions do not obstruct the NC implication: their output bits
can be computed in NC by the standard polynomial-configuration-graph
construction. A logspace machine has polynomially many configurations;
parallel Boolean matrix powering determines reachability with polynomial
size and $O(\log^2 n)$ depth. For a polynomial-output transducer, include a
logarithmic-size output-position counter and mark the transitions emitting
the requested bit. Adjacency tests are local transition checks with input
multiplexers, also uniformly shallow. Composing those circuits with an NC
circuit-value algorithm preserves polynomial size, polylogarithmic depth,
and logspace uniformity.

It follows that the missing general block-summary construction would solve
a P-complete target, not merely speed up one example. Its uniformity cannot
be delegated to an unrestricted offline optimizer: selecting arbitrary small
circuits separately at each length changes the class being studied.

\paragraph{Verification and outcome.}
The companion exact checks compare balanced and sequential composition for
all triples of functions on a three-element state set, and for all pairs
of affine maps on two bits. They exhaustively verify the carry-composition
identity, including redundant pair encodings, and test small binary additions.
The Fibonacci unfolding counts are checked separately from the constant-size
functional simplification. These checks validate the displayed algebra;
the resource bounds and uniformity rely on the explicit gate constructions.

\textbf{UNRESOLVED.} The attempt supplies fully specified parallel algorithms
for explicit finite-state updates, affine updates, and carry chains. It also
exhibits exact counterexamples to inferring hardness from sequential depth
or from an inefficient unfolding. The first remaining upper-bound gap is a
polynomial-size summary for arbitrary computation blocks, closed under a
uniform polylogarithmic-depth composition operation with polynomial total
work. A separation would instead require a lower bound against every such
uniform circuit family, beyond the particular representations examined here.
Neither requirement has been met.

\subsection{Sources}
\begin{itemize}
\item[S1]J. Heemstra, J. Martens, and A. Wijs, Evaluating Massively Parallel Algorithms for DFA Minimisation, Equivalence Checking and Inclusion Checking, Introduction (2026).
\url{https://arxiv.org/abs/2508.20735}.
\textit{Formulation source.}
\item[S2]Parallel Computation, Matt Williamson.
\url{https://community.wvu.edu/~krsubramani/courses/sp09/cc/lecnotes/parallel.pdf}.
\textit{Further reference; not a proof endorsement.}
\item[S3]Complexity of the Freezing Majority Rule with L-shaped Neighborhoods.
\url{https://arxiv.org/html/2509.16065v1}.
\textit{Further reference; not a proof endorsement.}
\item[E1]Richard E. Ladner,
\emph{The Circuit Value Problem Is Log Space Complete for P},
ACM SIGACT News 7(1) (1975), pp. 18--20.
\url{https://doi.org/10.1145/990518.990519}.
\url{https://homes.cs.washington.edu/~ladner/site/index.php/publications/}.
\textit{Primary attribution for circuit-value completeness; author publication
record consulted 27 September 2026. The reduction and its relevance to the
uniform NC target are explained in the attempt.}
\end{itemize}
\end{document}
